% =====================================================================
%  EXPERIMENTAL VALIDATION PLAN FOR THE BOUNDED PHASE SPACE
%  FRAMEWORK: FROM THEORY TO BENCHMARKED COMPOUND IDENTIFICATION
%
%  A detailed specification of experiments, code architecture,
%  and language extensions required to close the gap between
%  the BPS/S-Entropy theoretical corpus and empirically scored
%  compound identification.
% =====================================================================

\documentclass[11pt,a4paper]{article}

\usepackage[utf8]{inputenc}
\usepackage[T1]{fontenc}
\usepackage{lmodern}
\usepackage[a4paper,top=2.4cm,bottom=2.4cm,left=2.4cm,right=2.4cm]{geometry}
\usepackage{amsmath,amssymb,amsthm,mathtools}
\usepackage{bm}
\usepackage{booktabs}
\usepackage{array}
\usepackage{longtable}
\usepackage{enumitem}
\usepackage{microtype}
\usepackage{xcolor}
\usepackage{listings}
\usepackage{graphicx}
\usepackage[round,sort&compress]{natbib}
\usepackage[colorlinks=true,linkcolor=blue!45!black,
            citecolor=blue!45!black,urlcolor=blue!45!black]{hyperref}
\usepackage{cleveref}
\usepackage{fancyhdr}
\usepackage{titlesec}
\usepackage{caption}

\pagestyle{fancy}
\fancyhf{}
\fancyhead[L]{\small Experimental Validation Plan}
\fancyhead[R]{\small Sachikonye}
\fancyfoot[C]{\thepage}
\renewcommand{\headrulewidth}{0.4pt}

% ---- Theorem-like environments ----
\newtheoremstyle{planstyle}{\topsep}{\topsep}{\itshape}{}{\bfseries}{.}{.5em}{}
\theoremstyle{planstyle}
\newtheorem{theorem}{Theorem}[section]
\newtheorem{lemma}[theorem]{Lemma}
\newtheorem{proposition}[theorem]{Proposition}
\newtheorem{corollary}[theorem]{Corollary}
\theoremstyle{definition}
\newtheorem{definition}[theorem]{Definition}
\newtheorem{requirement}[theorem]{Requirement}
\newtheorem{experiment}[theorem]{Experiment}
\newtheorem{deliverable}[theorem]{Deliverable}
\theoremstyle{remark}
\newtheorem{remark}[theorem]{Remark}

\crefname{experiment}{Experiment}{Experiments}
\Crefname{experiment}{Experiment}{Experiments}
\crefname{requirement}{Requirement}{Requirements}
\Crefname{requirement}{Requirement}{Requirements}
\crefname{deliverable}{Deliverable}{Deliverables}
\Crefname{deliverable}{Deliverable}{Deliverables}

% ---- Listing style ----
\definecolor{sskw}{RGB}{30,64,175}
\definecolor{sscom}{RGB}{22,101,52}
\definecolor{ssstr}{RGB}{159,18,57}
\definecolor{ssbg}{RGB}{250,250,250}
\lstdefinelanguage{shapeshifter}{
  morekeywords={import,objective,instrument,phase,validate,target_list,source,
                network,precursor,foreach,in},
  sensitive=true,
  morecomment=[l]{//},
  morestring=[b]",
  keywordstyle=\color{sskw}\bfseries,
  commentstyle=\color{sscom}\itshape,
  stringstyle=\color{ssstr},
  basicstyle=\ttfamily\footnotesize,
  backgroundcolor=\color{ssbg},
  frame=single,framesep=4pt,xleftmargin=6pt,xrightmargin=4pt,
  breaklines=true,showstringspaces=false,columns=fullflexible,tabsize=2
}
\lstdefinelanguage{rust}{
  morekeywords={pub,fn,struct,impl,enum,let,mut,self,use,mod,return,if,else,
                for,in,match,Ok,Err,Result,Option,Some,None,Vec,f64,u32,u64,
                i32,usize,bool,true,false,String,where,trait},
  sensitive=true,
  morecomment=[l]{//},
  morecomment=[s]{/*}{*/},
  morestring=[b]",
  keywordstyle=\color{sskw}\bfseries,
  commentstyle=\color{sscom}\itshape,
  stringstyle=\color{ssstr},
  basicstyle=\ttfamily\footnotesize,
  backgroundcolor=\color{ssbg},
  frame=single,framesep=4pt,xleftmargin=6pt,xrightmargin=4pt,
  breaklines=true,showstringspaces=false,columns=fullflexible,tabsize=2
}
\lstset{language=shapeshifter}

% ---- Shorthand ----
\newcommand{\R}{\mathbb{R}}
\newcommand{\N}{\mathbb{N}}
\newcommand{\Z}{\mathbb{Z}}
\newcommand{\Sv}{\mathbf{S}}
\newcommand{\Sk}{S_{\mathrm{k}}}
\newcommand{\St}{S_{\mathrm{t}}}
\newcommand{\Se}{S_{\mathrm{e}}}
\newcommand{\Parts}{\mathcal{P}}
\newcommand{\calG}{\mathcal{G}}
\newcommand{\calV}{\mathcal{V}}
\newcommand{\calE}{\mathcal{E}}
\newcommand{\nP}{n_{\mathrm{P}}}
\newcommand{\estar}{e^{*}}
\newcommand{\Tmap}{\mathcal{T}}
\newcommand{\shk}[1]{\texttt{#1}}

\title{\textbf{Experimental Validation Plan for the Bounded Phase Space
Framework}\\[6pt]
\large From Theory to Benchmarked Compound Identification:\\
Experiments, Code Architecture, and Language Extensions}

\author{
Kundai Farai Sachikonye\\
\small Technical University of Munich\\
\small Chair of Proteomics and Bioanalytics\\
\small \texttt{kundai.sachikonye@wzw.tum.de}}

\date{\today}

\begin{document}
\maketitle

\begin{abstract}
\noindent
The Bounded Phase Space (BPS) framework comprises a corpus of
theoretical papers---partition-state graph search for tandem mass
spectrometry, complete ion trajectory tracing, the Shapeshifter DSL for
experiment specification, database-free peptide identification via
categorical fragmentation networks, causal knowledge graph construction
by circuit propagation, and attention-allocated federated retrieval---together
with a Python implementation covering mzML extraction, S-entropy
coordinate transformation, phase-lock detection, and vector embedding.

The theoretical results are internally consistent and validated against
narrow datasets (NIST AC\_CAC for small molecules, PRIDE PXD000001 for
peptides, UC Davis for metabolomics). No result has been tested
head-to-head against established methods on community benchmarks, and
the cross-domain universality claim remains empirically ungrounded.

This document specifies the experimental programme required to close
that gap. We define seven experiments ordered by priority, each with
its dataset, ground truth, success criterion, and the theoretical claim
it tests. We then specify the code changes required to execute them:
three new block types and one structural extension to the Shapeshifter
DSL; eleven new namespaces comprising fifty-five standard library
operations; a Rust core library of thirteen modules replacing
approximately 2,500 lines of Python with 800 lines of type-safe,
visualisation-free computation; and a third failure category (principled
refusal) extending the DSL's operational semantics. Every specification
is stated with the invariant it must preserve or is explicitly marked as
a design decision subject to revision.
\end{abstract}

\noindent\textbf{Keywords:} bounded phase space; S-entropy coordinates;
SEBD-MS; Shapeshifter DSL; CASMI benchmark; transition state validation;
causal knowledge graph; attention allocation; federated retrieval;
experimental design.

\tableofcontents
\clearpage

% =====================================================================
\part{Current State and the Gap}
% =====================================================================

\section{Theoretical Corpus}
\label{sec:corpus}

The framework as currently developed consists of six authored papers,
each contributing a distinct theoretical component, and a Python
implementation providing data extraction, coordinate transformation,
phase-lock detection, and vector embedding. \Cref{tab:corpus} summarises
the papers, their central claims, and the datasets against which each
has been validated.

\begin{table}[htbp]
\centering
\small
\caption{Theoretical corpus with validation status.}
\label{tab:corpus}
\begin{tabular}{@{}p{3.2cm}p{4.8cm}p{3.5cm}p{2.5cm}@{}}
\toprule
\textbf{Paper} & \textbf{Central claim} & \textbf{Validated against} &
\textbf{Comparison} \\
\midrule
Partition-State Graph (SEBD-MS) &
  MS/MS identification as shortest-path in partition-state graph; virtual substates
  are transition states &
  3,000 NIST AC\_CAC HCD spectra &
  None \\[4pt]
Ion Trajectory &
  All four analyzer types read the same partition count (Time-Count Identity) &
  Theoretical; single-ion walkthrough &
  None \\[4pt]
Shapeshifter DSL &
  Experiment specification as single executable document with proved determinism,
  isolation, monotonicity, stage separation &
  Formal proofs only; no executed programs &
  None \\[4pt]
Categorical Fragmentation (Lavoisier) &
  Database-free peptide identification via droplet encoding and fragment graphs;
  42\% partial match, 88.7\% PTM localisation, 89.3\% cross-platform transfer &
  PRIDE PXD000001 (1,247 spectra), 589 phosphopeptides, 400 multi-platform &
  MaxQuant \\[4pt]
Causal Knowledge Graphs (CKG) &
  CKG is propagation operator at rest; backward completion collapses from
  $\Theta(N)$ to $\Theta(\log_3 N)$ via virtual sub-states &
  Four metabolic networks; 17 consistency checks &
  None \\[4pt]
Attention Allocation &
  Join yield is sigmoid; threshold reverses optimal policy;
  allocator refuses below computable threshold &
  21 checks, 6 negative controls; closed-form validation &
  Brute-force exhaustive search \\
\bottomrule
\end{tabular}
\end{table}

\section{The Gap}
\label{sec:gap}

Three categories of gap separate the current state from publishable
experimental results.

\subsection{No community benchmark comparison}

No paper has been tested against established methods on a dataset the
community uses for comparison. SEBD-MS has not been run on CASMI.
The Lavoisier peptide pipeline has not been compared against SIRIUS,
MetFrag, or CFM-ID. The CKG has not been compared against curated
databases on a shared network.

\subsection{No cross-domain validation}

Each paper validates against one dataset in one domain. The
universality claim---that partition states and S-entropy coordinates
apply across small molecules, peptides, metabolites, and metabolic
networks---has no empirical support spanning two of these domains
simultaneously.

\subsection{No executable specification}

The Shapeshifter DSL has formal proofs but no executed programs against
real data. The Python implementation is imperative code that the DSL
cannot currently invoke, because the DSL lacks the block types, standard
library operations, and structural extensions required to express the
experiments described here.

% =====================================================================
\part{Experimental Programme}
% =====================================================================

\section{Design Principles}
\label{sec:principles}

Each experiment is specified by five components: the \emph{dataset}
(what data is used), the \emph{ground truth} (what the correct answer
is and how it was determined), the \emph{theoretical claim} (which
theorem or proposition is being tested), the \emph{success criterion}
(a quantitative threshold stated before results are seen), and the
\emph{comparison method} (an established tool or published result
against which the framework's output is scored). An experiment without
a comparison method is a characterisation study, not a validation.

\section{Experiment 1: CASMI 2022 Benchmark}
\label{sec:exp-casmi}

\begin{experiment}[CASMI 2022 Compound Identification]
\label{exp:casmi}
\textbf{Objective.} Score SEBD-MS and ternary trie lookup against the
CASMI 2022 challenge set on all four evaluation categories.

\textbf{Dataset.} UC Davis CASMI 2022 challenge: 500 compounds (250
priority, 250 bonus) with raw LC-MS/MS data in positive and negative ESI,
provided as mzML with retention times and precursor $m/z$ values. Solutions
published at the Google Drive repository maintained by the West Coast
Metabolomics Center.

\textbf{Ground truth.} Published solutions comprising InChIKey, molecular
formula, SMILES, and compound class for each challenge compound. Ground
truth was determined by the organisers using authentic standards.

\textbf{Theoretical claims tested.}
\begin{enumerate}[nosep,label=(\roman*)]
\item SEBD-MS (Theorem~4.2 of the partition-state graph paper) finds
  minimum-cost paths corresponding to correct identifications.
\item The ternary trie provides $\mathcal{O}(k_0)$ lookup independent
  of database size (Corollary~7.2).
\item CKG-derived constraints improve identification over unconstrained
  search (Theorem~4.5 of the CKG paper).
\item The attention allocator correctly distributes query budget across
  federated sources (Theorem~3.1 of the allocation paper).
\end{enumerate}

\textbf{Success criteria.}
\begin{itemize}[nosep]
\item Category 1 (adduct annotation): $\geq 70\%$ correct.
\item Category 2 (molecular formula): $\geq 50\%$ correct.
\item Category 3 (compound class): $\geq 40\%$ correct.
\item Category 4 (2D structure, InChIKey first block): $\geq 20\%$
  rank-1 accuracy.
\end{itemize}

These thresholds are set conservatively below published CASMI 2022
results from SIRIUS+CSI:FingerID and MetFrag, so that meeting them
establishes viability rather than claiming superiority.

\textbf{Comparison methods.} Published CASMI 2022 submissions:
SIRIUS+CSI:FingerID, MetFrag, CFM-ID, MS-FINDER. Scores are available
from the workshop proceedings.

\textbf{Protocol.}
\begin{enumerate}[nosep]
\item Extract spectra from mzML files using \shk{lavoisier.extract.load}.
\item Build ternary tries over HMDB (220K), MassBank (93K), and
  GNPS (500K) reference libraries.
\item Allocate query budget across the three sources using
  \shk{lavoisier.allocate.gated} with per-source thresholds computed
  from corpus sizes.
\item For each challenge spectrum: compute S-entropy coordinates,
  run SEBD-MS backward and forward search, query each trie with
  allocated budget, and collect candidate lists.
\item Optionally: solve the central metabolism circuit, recover the CKG,
  compute the reachable set from already-identified compounds, and
  re-rank candidates using the CKG constraint.
\item Score against published solutions on all four categories.
\item Report: rank-1 accuracy, top-5 accuracy, computation time per
  spectrum, and the CKG-constrained versus unconstrained comparison.
\end{enumerate}
\end{experiment}

\section{Experiment 2: Transition State Geometry}
\label{sec:exp-ts}

\begin{experiment}[Virtual Substates vs.\ DFT Transition States]
\label{exp:ts}
\textbf{Objective.} Test whether off-shell virtual substates found by
SEBD-MS backward search correspond to DFT-computed transition state
geometries for known fragmentation mechanisms.

\textbf{Dataset.} Published MS/MS spectra of compounds with
well-characterised fragmentation mechanisms: McLafferty rearrangement
(ketones, esters), retro-Diels-Alder (cyclohexenes), $\alpha$-cleavage
(amines, ethers). DFT saddle-point geometries from published
computational chemistry literature.

\textbf{Ground truth.} DFT-optimised transition state structures with
published energies and geometries (bond lengths, angles, dihedrals) at
the B3LYP/6-31G(d) or M06-2X/6-311+G(2d,2p) level.

\textbf{Theoretical claim tested.} Theorem~5.4 of the partition-state
graph paper: every off-shell component $\Sv_i \notin [0,1]^3$ in an
optimal virtual decomposition of a fragment node is a transition state
in the fragmentation pathway.

\textbf{Success criterion.} Statistically significant correlation
($p < 0.01$) between off-shell S-entropy coordinates and DFT
saddle-point geometry projected into S-entropy space, for at least
5 of 10 tested reaction mechanisms.

\textbf{Comparison methods.} No existing tool connects graph-search
intermediates to quantum-chemical transition states. A positive result
here is novel with no prior comparison.

\textbf{Protocol.}
\begin{enumerate}[nosep]
\item Select 10 compounds with published fragmentation mechanisms and
  DFT transition state geometries.
\item For each compound: obtain experimental MS/MS spectrum (from NIST
  or MassBank), run SEBD-MS with full backward search, extract all
  virtual substate coordinates via
  \shk{lavoisier.msms.extract\_virtual\_substates}.
\item Project DFT transition state geometry into S-entropy space using
  the partition bijection $\varPhi$ and coordinate embedding.
\item Compute Euclidean distance between each virtual substate and the
  projected DFT geometry.
\item Report: correlation coefficient, per-mechanism agreement, and
  the distribution of virtual-to-DFT distances.
\end{enumerate}
\end{experiment}

\section{Experiment 3: Cross-Platform Lagrangian Convergence}
\label{sec:exp-lagrangian}

\begin{experiment}[Time-Count Identity Across Analyzers]
\label{exp:lagrangian}
\textbf{Objective.} Test whether S-entropy coordinates computed
independently from different mass analyzer platforms converge for the
same compounds.

\textbf{Dataset.} NIST SRM 1950 (Metabolites in Frozen Human Plasma),
measured on Orbitrap, FT-ICR, TOF, and triple-quad instruments by
independent laboratories. Approximately 100 metabolites with
cross-platform measurements available.

\textbf{Ground truth.} Compound identities established by consensus
across laboratories; mass accuracy certified by NIST.

\textbf{Theoretical claim tested.} Time-Count Identity (Theorem~6.2 of
the ion trajectory paper): $d\mathcal{M}/dt = \omega/(2\pi) = 1/\langle\tau_p\rangle$
exactly for every analyzer, so S-entropy coordinates from any analyzer
should give the same embedding.

\textbf{Success criterion.} Maximum pairwise $L_2$ distance between
S-entropy coordinates computed from different platforms for the same
compound: $d_{\max} < 0.05$ for $\geq 80\%$ of compounds.

\textbf{Comparison methods.} No existing framework predicts
platform-convergence of spectral coordinates. The Lavoisier paper's
cross-platform CV $< 2.1\%$ for topology features is the internal
comparison; this experiment tests the underlying coordinate convergence
that topology invariance depends on.

\textbf{Protocol.}
\begin{enumerate}[nosep]
\item Obtain SRM 1950 spectra from each platform.
\item For each compound on each platform: extract spectrum, compute
  S-entropy coordinates via \shk{lavoisier.sentropy.transform}, compute
  partition coordinates via \shk{lavoisier.observe.partition\_from\_sentropy}.
\item Optionally: compute circuit-derived coordinates from a metabolic
  network context via \shk{lavoisier.circuit.node\_coordinates}, providing
  a third independent derivation.
\item For each compound: compute pairwise $L_2$ distances across platforms.
\item Report: distribution of pairwise distances, maximum deviation,
  and the fraction of compounds meeting the convergence threshold.
  Systematic platform-dependent offsets, if any, are reported as
  diagnostic rather than averaged away.
\end{enumerate}
\end{experiment}

\section{Experiment 4: Glycopeptide Fragmentation}
\label{sec:exp-glyco}

\begin{experiment}[Forward Reachability on Nonlinear Structures]
\label{exp:glyco}
\textbf{Objective.} Test whether the SEBD-MS forward reachability cone
correctly contains observed glycan fragments from glycopeptide
precursors, stress-testing the framework on branching (nonlinear)
fragmentation.

\textbf{Dataset.} Published glycopeptide HCD spectra with known
structures from the Struwe and Wuhrer groups. Glycan fragments include
oxonium ions (HexNAc 204.087, Hex 163.060, NeuAc 292.103) and Y-ions
from glycosidic cleavage.

\textbf{Ground truth.} Glycan composition and peptide sequence
determined by orthogonal methods (enzymatic release, HPAEC-PAD, or
synthetic standards).

\textbf{Theoretical claim tested.} The forward reachability cone
$\mathcal{R}_f^{(d)}$ from the precursor (Definition~6.1 of the
partition-state graph paper) should contain all observed glycan
fragments if the partition-state graph correctly handles branching
structures.

\textbf{Success criterion.} $\geq 85\%$ of observed glycan-specific
fragments fall within $\mathcal{R}_f^{(d_{\max})}$ for $d_{\max} = 10$.

\textbf{Protocol.}
\begin{enumerate}[nosep]
\item For each glycopeptide: compute the precursor's partition state,
  build the forward reachability cone to depth 10 via
  \shk{lavoisier.msms.reachability\_cone}.
\item Check membership of each observed fragment via
  \shk{lavoisier.msms.cone\_membership}.
\item Report: membership rate for glycan-specific fragments versus
  peptide backbone fragments; identify systematic failures, if any,
  attributable to branching topology.
\end{enumerate}
\end{experiment}

\section{Experiment 5: MassBank-Scale Ternary Trie Scaling}
\label{sec:exp-scaling}

\begin{experiment}[Database-Size Independence of Query Cost]
\label{exp:scaling}
\textbf{Objective.} Verify that SEBD-MS per-query cost is independent
of database size, as claimed by Corollary~7.2 of the partition-state
graph paper.

\textbf{Dataset.} MassBank (93K spectra), GNPS MassIVE subset (500K),
and optionally PubChem-scale candidate lists ($10^8$). Query set: the
CASMI 2022 challenge spectra from \cref{exp:casmi}.

\textbf{Ground truth.} Wall-clock time is the measurement; the
theoretical prediction is $\mathcal{O}(k_0)$ per query, independent
of $N_{\text{db}}$.

\textbf{Theoretical claim tested.} Corollary~7.2: per-query cost is
$\mathcal{O}(k_0 \cdot n_p d_{\max} \log(n_p d_{\max}))$, independent
of $N_{\text{db}}$. Additionally, Corollary~5.7: empty lookups at
on-shell addresses predict unknown compounds.

\textbf{Success criterion.} Wall-clock time per query varies by $< 20\%$
across a 100-fold range of database sizes ($N_{\text{db}}$ from 1K to
100K).

\textbf{Protocol.}
\begin{enumerate}[nosep]
\item Build ternary tries at depths $k_0 \in \{9, 12, 15, 18\}$ over
  subsampled databases of sizes $\{1\text{K}, 5\text{K}, 10\text{K},
  50\text{K}, 93\text{K}\}$.
\item Run a fixed set of 100 query spectra against each trie.
\item Record wall-clock time per query via
  \shk{validate.timing\_profile}.
\item Plot query time versus database size; fit slope.
\item Report: coefficient of variation of query time across database
  sizes; empty-lookup addresses and whether any correspond to compounds
  added to MassBank in later versions.
\end{enumerate}
\end{experiment}

\section{Experiment 6: Federated Retrieval with Attention Allocation}
\label{sec:exp-federated}

\begin{experiment}[Gated Allocation Across Federated Sources]
\label{exp:federated}
\textbf{Objective.} Test the gated attention allocator on a real
federated retrieval task, verifying that normalised efforts equalise
and that the allocator refuses below threshold.

\textbf{Dataset.} Three federated spectral databases: MassBank (93K),
HMDB (220K), GNPS (500K). Query budget: variable from $0.1\estar$ to
$5\estar$ of the smallest source.

\textbf{Ground truth.} The allocator's output is compared against
brute-force exhaustive search over the allocation simplex, which shares
no reasoning with the allocator.

\textbf{Theoretical claims tested.}
\begin{enumerate}[nosep,label=(\roman*)]
\item Theorem~2.1 (allocation paper): $\estar = 2\ln 2 / \ln(1/(1-1/N))$
  is exact and $Y(\estar) = N/4$.
\item Theorem~3.1: below $\estar$ committing beats splitting.
\item Proposition~3.3: the optimum equalises $e/\estar$, not $e$.
\item Definition~4.2: the allocator refuses below $\min_i \estar_i$ and
  reports the required budget.
\end{enumerate}

\textbf{Success criterion.} Allocator yield within $10^{-4}$ of
brute-force yield at every tested budget. Refusal fires below threshold
and does not fire above.

\textbf{Protocol.}
\begin{enumerate}[nosep]
\item Compute $\estar_i$ for each source from corpus sizes.
\item Sweep budget from $0.1 \cdot \min_i \estar_i$ to $5 \cdot \sum_i \estar_i$.
\item At each budget: run \shk{lavoisier.allocate.gated}, run
  brute-force search, compare yields.
\item Report: yield gap, normalised efforts, refusal points, and
  the budget at which the refusal resolves.
\end{enumerate}
\end{experiment}

\section{Experiment 7: Shapeshifter DSL Validation}
\label{sec:exp-dsl}

\begin{experiment}[Specification Fidelity]
\label{exp:dsl}
\textbf{Objective.} Demonstrate that Shapeshifter programs faithfully
encode published experimental methods and that execution reproduces
published results.

\textbf{Dataset.} Five published methods sections from metabolomics and
proteomics papers with fully specified protocols.

\textbf{Ground truth.} The published results (compound lists, retention
times, mass accuracies) from each paper.

\textbf{Theoretical claim tested.} The specification-gap thesis of the
Shapeshifter paper: a single executable document captures what a methods
section describes.

\textbf{Success criterion.} For each of five methods sections, the
\shk{.ss} program compiles without error and produces output matching
the published results to within stated mass accuracy and retention time
tolerance.

\textbf{Protocol.}
\begin{enumerate}[nosep]
\item Select five published papers with complete methods sections.
\item Encode each as a \shk{.ss} program using the extended grammar.
\item Compile each program (\textsc{Compile} stage, no execution).
\item Execute each program (\textsc{Execute} stage, against available data).
\item Compare output against published results.
\end{enumerate}
\end{experiment}

\section{Experiment Dependency Structure}
\label{sec:dependencies}

The experiments have a partial order dictated by shared infrastructure.
\Cref{tab:dependencies} records the ordering.

\begin{table}[htbp]
\centering
\small
\caption{Experiment dependencies. An entry $(i, j)$ means experiment $i$
requires infrastructure built for experiment $j$.}
\label{tab:dependencies}
\begin{tabular}{@{}lccccccc@{}}
\toprule
 & Exp.~1 & Exp.~2 & Exp.~3 & Exp.~4 & Exp.~5 & Exp.~6 & Exp.~7 \\
\midrule
Exp.~1 (CASMI)       & ---  &      &      &      & $\leftarrow$ & $\leftarrow$ & \\
Exp.~2 (Trans.\ St.) &      & ---  &      &      &       &       & \\
Exp.~3 (Lagrangian)  &      &      & ---  &      &       &       & \\
Exp.~4 (Glyco)       &      &      &      & ---  &       &       & \\
Exp.~5 (Scaling)     &      &      &      &      & ---   &       & \\
Exp.~6 (Federated)   &      &      &      &      &       & ---   & \\
Exp.~7 (DSL)         & $\leftarrow$ & & & & & $\leftarrow$ & --- \\
\bottomrule
\end{tabular}
\end{table}

Experiment~2 (transition states) and Experiment~3 (Lagrangian
convergence) are independent of all others. Experiment~1 (CASMI)
depends on the trie infrastructure of Experiment~5 and the allocator
of Experiment~6. Experiment~7 (DSL) depends on Experiments~1 and~6
for the programs it encodes.

The recommended execution order for maximum information per unit of
implementation effort is:

\begin{enumerate}[nosep]
\item Experiment~2 (transition states): highest novelty, requires no
  new data collection, tests the most distinctive theoretical claim.
\item Experiment~5 (trie scaling): required infrastructure for
  Experiment~1; pure algorithmic validation.
\item Experiment~6 (federated allocation): required infrastructure for
  Experiment~1; closed-form validation.
\item Experiment~1 (CASMI): the benchmark that gives everything else
  credibility.
\item Experiment~3 (Lagrangian convergence): independent, uses existing
  SRM 1950 data.
\item Experiment~4 (glycopeptide): stress-test on nonlinear structures.
\item Experiment~7 (DSL validation): requires most infrastructure,
  provides software-engineering rather than scientific evidence.
\end{enumerate}

% =====================================================================
\part{Language Extensions}
% =====================================================================

\section{Design Constraints}
\label{sec:design-constraints}

All extensions to the Shapeshifter DSL must preserve the four proved
properties of the original specification: declaration-order determinism
(Theorem~5.3), phase isolation (Theorem~5.4), workspace monotonicity
(Theorem~5.5), and stage separation (Theorem~5.7). Each extension is
stated with the invariant it preserves or modifies.

\section{New Block Types}
\label{sec:blocks}

\subsection{The \texttt{source} block}

\begin{requirement}[\texttt{source} block]
\label{req:source}
A \shk{source} block is a field block (no executable statements)
declaring external data to be consumed by standard library operations.
Grammar:
\begin{lstlisting}
source  ::=  "source" IDENT ":" fieldblock
\end{lstlisting}
Required fields: \shk{format} (one of \shk{"mzml"}, \shk{"mgf"},
\shk{"msp"}, \shk{"csv"}), \shk{path} (string). Optional fields:
\shk{vendor}, \shk{provenance} (free text), \shk{ground\_truth} (path),
\shk{ground\_truth\_key} (field name in the truth file),
\shk{instrument} (bare-word reference to an \shk{instrument} block).

\emph{Invariant:} \shk{source} is a field block, so Theorems~5.3--5.7
hold without modification. The compile stage validates that \shk{format}
and \shk{path} are present; it does not check file existence (stage
separation preserved).
\end{requirement}

\subsection{The \texttt{network} block}

\begin{requirement}[\texttt{network} block]
\label{req:network}
A \shk{network} block declares a reaction network for CKG construction.
Grammar:
\begin{lstlisting}
network  ::=  "network" IDENT ":" fieldblock
\end{lstlisting}
Required fields: \shk{species} (array of objects with \shk{name},
\shk{mu\_std}, \shk{concentration}), \shk{reactions} (array of objects
with \shk{from}, \shk{to}, \shk{k\_cat}), \shk{temperature} (number
in Kelvin).

\emph{Invariant:} Same as \shk{source}---field block, no execution,
all theorems preserved.
\end{requirement}

\subsection{The \texttt{precursor} block}

\begin{requirement}[\texttt{precursor} block]
\label{req:precursor}
A \shk{precursor} block declares a target ion for structured
precursor-fragment operations.
Grammar:
\begin{lstlisting}
precursor  ::=  "precursor" IDENT ":" fieldblock
\end{lstlisting}
Required fields: \shk{mz} (number), \shk{charge} (integer). Optional:
\shk{adduct} (string), \shk{rt\_range} (pair), \shk{tolerance\_ppm}.

\emph{Invariant:} Field block; theorems preserved.
\end{requirement}

\section{Structural Extension: The \texttt{foreach} Modifier}
\label{sec:foreach}

\begin{requirement}[\texttt{foreach} phase modifier]
\label{req:foreach}
Grammar:
\begin{lstlisting}
phase  ::=  "phase" IDENT ["foreach" IDENT "in" IDENT] ":" stmtblock
\end{lstlisting}
Semantics: the body is a template instantiated once per element of the
named source or list, in source order. Bindings produced across
iterations are collected into a list under the original name.

\emph{Invariants preserved:}
\begin{enumerate}[nosep,label=(\roman*)]
\item \emph{Declaration-order determinism} (Theorem~5.3): elements are
  processed in source order; the body is fixed.
\item \emph{Workspace monotonicity} (Theorem~5.5): the collected list
  is appended once after all iterations complete.
\item \emph{Enumerability} (Proposition~5.9): call heads in the body
  are visible by inspection.
\end{enumerate}

\emph{Invariant modified:} Phase isolation (Theorem~5.4) gains a
scoped exception: the iteration variable is local to each instantiation
of the body. This is the single exception to the ``no phase-local
binding'' rule and is stated explicitly in the modified theorem.
\end{requirement}

\section{Failure Classification Extension: Principled Refusal}
\label{sec:refusal}

\begin{requirement}[Refusal as a third failure category]
\label{req:refusal}
Proposition~5.8 of the original specification defines two failure
categories: local failure (warning) and global failure (error). We
add a third:

\textbf{(iii) Principled refusal.} A condition under which the
operation has a well-defined answer that is not a result but a
diagnostic statement about why no result is defensible. The binding
receives kind \shk{object} with a distinguished field
\shk{refused: true}. The refusal carries: the offered value, the
required value, the shortfall, and per-component diagnostics.

The refusal does not halt execution (unlike an error) and cannot be
consumed as a result by downstream operations (unlike a warning-tagged
binding). It is appended to the log with level \shk{refusal}.

\emph{Source:} Definition~4.2 and Principle~4.1 of the attention
allocation paper; Definition~4.2 of the medium-vertex paper.
\end{requirement}

% =====================================================================
\part{Standard Library Extensions}
% =====================================================================

\section{Overview}
\label{sec:stdlib-overview}

The extended standard library comprises eleven namespaces and fifty-five
operations. \Cref{tab:stdlib-summary} gives the counts; the full
specification of each operation follows in
\crefrange{sec:ns-extract}{sec:ns-validate}.

\begin{table}[htbp]
\centering
\small
\caption{Standard library namespace summary.}
\label{tab:stdlib-summary}
\begin{tabular}{@{}llcc@{}}
\toprule
\textbf{Namespace} & \textbf{Domain} & \textbf{Operations} &
\textbf{New} \\
\midrule
\shk{lavoisier.instrument}  & Experiment generation     & 2 & 0 \\
\shk{lavoisier.extract}     & mzML/spectral data I/O    & 6 & 6 \\
\shk{lavoisier.sentropy}    & S-entropy transformation  & 4 & 4 \\
\shk{lavoisier.partition}   & Partition coordinates     & 1 & 0 \\
\shk{lavoisier.phaselock}   & Phase-lock detection      & 7 & 7 \\
\shk{lavoisier.vectors}     & Vector embeddings         & 4 & 4 \\
\shk{lavoisier.msms}        & SEBD-MS and fragmentation & 10 & 4 \\
\shk{lavoisier.db}          & Database/trie operations  & 6 & 3 \\
\shk{lavoisier.circuit}     & CKG construction          & 7 & 7 \\
\shk{lavoisier.allocate}    & Attention allocation      & 4 & 4 \\
\shk{lavoisier.compare}     & Cross-pipeline comparison & 4 & 4 \\
\shk{validate}              & Validation operations     & 6 & 5 \\
\midrule
\textbf{Total}              &                           & \textbf{61} & \textbf{48} \\
\bottomrule
\end{tabular}
\end{table}

\section{Namespace: \texttt{lavoisier.extract}}
\label{sec:ns-extract}

These operations replace the imperative Python functions
\shk{extract\_mzml}, \shk{MSDataContainer}, \shk{get\_spectra}, and
\shk{get\_xic\_from\_pl} with declarative operations whose parameters
are exposed as named arguments.

\begin{longtable}{@{}p{3.8cm}p{5.5cm}p{1.8cm}p{1.5cm}@{}}
\toprule
\textbf{Operation} & \textbf{Arguments} & \textbf{Kind} & \textbf{Totality} \\
\midrule
\endhead
\shk{load} &
  \shk{source}, \shk{instrument}, \shk{rt\_range}, \shk{dda\_top?}$=6$,
  \shk{ms1\_threshold?}$=1000$, \shk{ms2\_threshold?}$=10$,
  \shk{ms1\_max?}$=0$, \shk{min\_spec\_peaks?}$=3$ &
  \shk{records} & partial \\
\addlinespace
\shk{organize} & \shk{raw\_data} & \shk{records} & total \\
\shk{link\_precursors} & \shk{organized}, \shk{precision?} & \shk{records} & total \\
\shk{xic} & \shk{source}, \shk{target\_mz}, \shk{tolerance\_ppm?}$=10$,
  \shk{rt\_range?} & \shk{records} & total \\
\shk{xic\_batch} & \shk{source}, \shk{targets} (list), \shk{tolerance\_ppm?},
  \shk{rt\_range?} & \shk{records} & total \\
\shk{metadata} & \shk{source} & \shk{object} & total \\
\bottomrule
\end{longtable}

The \shk{vendor} field resides on the \shk{source} block rather than
the \shk{instrument} block, because vendor-specific parsing logic
affects how raw data is structured (e.g., Waters multi-function DDA),
not the physics of the instrument.

\section{Namespace: \texttt{lavoisier.sentropy}}
\label{sec:ns-sentropy}

These operations replace \shk{SEntropyTransformer} with declarative
operations exposing the four configurable parameters
($\alpha$, $\beta$, $k$-neighbors, minimum intensity).

\begin{longtable}{@{}p{3.8cm}p{5.5cm}p{1.8cm}p{1.5cm}@{}}
\toprule
\textbf{Operation} & \textbf{Arguments} & \textbf{Kind} & \textbf{Totality} \\
\midrule
\endhead
\shk{transform} & \shk{spectra}, \shk{alpha?}$=0.5$, \shk{beta?}$=1.0$,
  \shk{k\_neighbors?}$=5$, \shk{min\_intensity?}$=0$ &
  \shk{records} & total \\
\shk{extract\_features} & \shk{coordinates} & \shk{records} & total \\
\shk{transform\_and\_extract} & (combined) & \shk{records} & total \\
\shk{pairwise\_distances} & \shk{coordinates} & \shk{tensor} & total \\
\shk{from\_retention} & \shk{retention\_times}, \shk{column}, \shk{k\_factor?}$=1$ &
  \shk{records} & total \\
\shk{to\_retention} & \shk{coordinates}, \shk{column}, \shk{k\_factor?}$=1$ &
  \shk{records} & total \\
\bottomrule
\end{longtable}

The \shk{from\_retention} and \shk{to\_retention} operations replace
\shk{PartitionLagCalculator} from the chromatographic trap array module.

\section{Namespace: \texttt{lavoisier.circuit}}
\label{sec:ns-circuit}

\begin{longtable}{@{}p{3.8cm}p{5.5cm}p{1.8cm}p{1.5cm}@{}}
\toprule
\textbf{Operation} & \textbf{Arguments} & \textbf{Kind} & \textbf{Totality} \\
\midrule
\endhead
\shk{solve} & \shk{network} & \shk{object} & partial \\
\shk{propagate} & \shk{circuit}, \shk{process?}$=$\shk{"rest"} &
  \shk{records} & total \\
\shk{reachable\_set} & \shk{ckg}, \shk{identified}, \shk{depth?}$=3$ &
  \shk{records} & total \\
\shk{contact\_map} & \shk{circuit} & \shk{tensor} & total \\
\shk{node\_coordinates} & \shk{circuit}, \shk{species} & \shk{object} & total \\
\shk{coherence} & \shk{circuit} & \shk{scalar} & total \\
\shk{axis\_reducibility} & \shk{circuit} & \shk{object} & total \\
\bottomrule
\end{longtable}

The \shk{solve} operation is partial because it requires the Wegscheider
condition (Theorem~2.2 of the CKG paper). Violation is a global failure:
without a consistent potential function, no circuit and no CKG exist.

\section{Namespace: \texttt{lavoisier.allocate}}
\label{sec:ns-allocate}

\begin{longtable}{@{}p{3.8cm}p{5.5cm}p{1.8cm}p{1.5cm}@{}}
\toprule
\textbf{Operation} & \textbf{Arguments} & \textbf{Kind} & \textbf{Totality} \\
\midrule
\endhead
\shk{gated} & \shk{sources}, \shk{budget},
  \shk{threshold\_method?}$=$\shk{"quarter\_corpus"} &
  \shk{object} & partial (may refuse) \\
\shk{threshold} & \shk{corpus\_size} & \shk{scalar} & total \\
\shk{quarter\_corpus\_check} & \shk{yield\_count}, \shk{corpus\_size} &
  \shk{scalar} & total \\
\shk{normalised\_effort} & \shk{effort}, \shk{corpus\_size} &
  \shk{scalar} & total \\
\bottomrule
\end{longtable}

The \shk{gated} operation is the only operation in the standard library
that may produce a principled refusal (\cref{req:refusal}).

\section{Namespace: \texttt{validate} (extended)}
\label{sec:ns-validate}

\begin{longtable}{@{}p{3.8cm}p{5.5cm}p{1.8cm}p{1.5cm}@{}}
\toprule
\textbf{Operation} & \textbf{Arguments} & \textbf{Kind} & \textbf{Totality} \\
\midrule
\endhead
\shk{score\_against\_truth} & \shk{results}, \shk{source}, \shk{metric}
  (\shk{"rank1\_accuracy"}, \shk{"top5\_accuracy"}, \shk{"formula\_accuracy"},
  \shk{"class\_accuracy"}) & \shk{scalar} & partial \\
\shk{compare\_results} & \shk{result\_a}, \shk{result\_b}, \shk{key},
  \shk{metric} & \shk{object} & total \\
\shk{assert\_convergence} & \shk{results} (list), \shk{metric},
  \shk{threshold} & \shk{object} & total \\
\shk{assert\_within} & \shk{value}, \shk{bounds} & \shk{scalar} & total \\
\shk{timing\_profile} & \shk{results}, \shk{label} & \shk{object} & total \\
\shk{check\_resolution\_time} & \shk{instrument}, \shk{window\_ppm} &
  (existing) & total \\
\bottomrule
\end{longtable}

The \shk{score\_against\_truth} operation is partial because it requires
the \shk{ground\_truth} field on the \shk{source} block. If
\shk{ground\_truth} is absent, the operation appends a warning (not an
error) and returns $-1$, consistent with the failure classification:
the experiment is meaningful without scoring, but the score is not
available.

% =====================================================================
\part{Rust Core Library}
% =====================================================================

\section{Architecture}
\label{sec:rust-arch}

The Rust core library provides the computational backend for all
standard library operations. It replaces approximately 2,500 lines of
Python across five modules (\shk{DataContainer}, \shk{EntropyTransformation},
\shk{PhaseLockDetectionNetworks}, \shk{VectorTransformation},
\shk{ChromatographicTrapArray}) with approximately 800 lines of
type-safe Rust distributed across thirteen modules. The Rust modules
contain no I/O, no visualisation, and no class hierarchies.

\begin{table}[htbp]
\centering
\small
\caption{Rust module structure. Lines are estimates based on the
specified interfaces.}
\label{tab:rust-modules}
\begin{tabular}{@{}llll@{}}
\toprule
\textbf{Module} & \textbf{Contains} & \textbf{Source} & \textbf{Est.\ lines} \\
\midrule
\shk{constants.rs}     & Physical constants              & All papers   & 15 \\
\shk{sentropy.rs}      & \shk{SEntropy}, magnitude, distance,
                         from/to retention                & Papers 1,2; Doc 11 & 80 \\
\shk{partition.rs}     & \shk{PartitionState}, bijection index,
                         capacity, from\_sentropy         & Paper 1; Doc 11 & 90 \\
\shk{ternary.rs}       & \shk{TernaryAddress}, encode, decode,
                         navigate, common\_prefix         & Paper 1; CKG & 70 \\
\shk{trie.rs}          & Lazy dictionary, empty lookup    & Paper 1      & 100 \\
\shk{column.rs}        & \shk{ColumnParams}, RT $\leftrightarrow$
                         S-entropy, undetermined residue  & Doc 11       & 60 \\
\shk{analyzer.rs}      & Cyclotron freq, TOF time,
                         Orbitrap freq, quad stability    & Paper 2      & 50 \\
\shk{sebd.rs}          & SEBD-MS algorithm, forward/backward
                         search, virtual substates        & Paper 1      & 120 \\
\shk{graph.rs}         & Fragment graph, admissible transitions,
                         Hamiltonian path                 & Paper 4      & 80 \\
\shk{circuit.rs}       & \shk{ReactionNetwork}, solve,
                         Kirchhoff, Wegscheider           & CKG          & 90 \\
\shk{propagation.rs}   & $\Tmap$ operator, CKG recovery,
                         reachable set                    & CKG          & 70 \\
\shk{allocator.rs}     & Gated allocation, refusal,
                         threshold, normalised effort     & Allocation   & 80 \\
\shk{equivalence.rs}   & \shk{TripleEquivalence},
                         agreement checks                 & Doc 12       & 40 \\
\shk{validation.rs}    & Resolution, distinguishability,
                         panel floor                      & CKG; Doc 11  & 50 \\
\midrule
\textbf{Total}         &                                  &              & $\sim$\textbf{895} \\
\bottomrule
\end{tabular}
\end{table}

\section{Core Types}
\label{sec:core-types}

Three types are shared across all modules and constitute the framework's
type-level invariants.

\begin{lstlisting}[language=rust]
/// S-Entropy coordinates in [0,1]^3.
/// Invariant: all components in [0,1] enforced at construction.
#[derive(Debug, Clone, Copy, PartialEq)]
pub struct SEntropy {
    pub s_k: f64,  // knowledge / throughput
    pub s_t: f64,  // time / capacity
    pub s_e: f64,  // entropy / potential
}

/// Partition state (n, l, m, s).
/// Invariant: n >= 1, 0 <= l < n, -l <= m <= l, s in {Up, Down}.
#[derive(Debug, Clone, Copy, PartialEq, Eq, Hash)]
pub struct PartitionState {
    pub n: u32,
    pub l: u32,
    pub m: i32,
    pub s: Spin,
}

/// Ternary address of depth k.
/// Invariant: each trit in {0, 1, 2}.
#[derive(Debug, Clone, PartialEq, Eq, Hash)]
pub struct TernaryAddress {
    pub trits: Vec<u8>,
}
\end{lstlisting}

\section{Module Interfaces}
\label{sec:interfaces}

We specify the public interface of each module by its function
signatures. Internal implementation is not specified; the interfaces
are the contract the DSL runtime depends on.

\subsection{\texttt{circuit.rs}}

\begin{lstlisting}[language=rust]
pub struct ReactionNetwork {
    pub species: Vec<Species>,
    pub reactions: Vec<Reaction>,
}

pub struct SolvedCircuit {
    pub potentials: Vec<f64>,      // mu_i
    pub conductances: Vec<f64>,    // G_ij
    pub fluxes: Vec<f64>,          // J_ij
    pub coordinates: Vec<SEntropy>,
    pub kirchhoff_residual: f64,
    pub wegscheider_holds: bool,
}

impl ReactionNetwork {
    pub fn solve(&self, temperature_k: f64)
        -> Result<SolvedCircuit, CircuitError>;
}

impl SolvedCircuit {
    pub fn contact_map(&self) -> Vec<(usize, usize, f64)>;
    pub fn coherence(&self) -> f64;
    pub fn axis_reducibility(&self) -> [f64; 3]; // R^2 per axis
}
\end{lstlisting}

\subsection{\texttt{allocator.rs}}

\begin{lstlisting}[language=rust]
pub struct JoinSource {
    pub name: String,
    pub corpus_size: u64,
}

pub enum AllocationResult {
    Allocated {
        per_source: Vec<(String, f64)>,
        normalised: Vec<(String, f64)>,
        policy: AllocationPolicy,
    },
    Refused {
        offered_budget: f64,
        required_budget: f64,
        shortfall: f64,
        per_source_thresholds: Vec<(String, f64)>,
    },
}

pub fn allocate(sources: &[JoinSource], budget: f64)
    -> AllocationResult;
pub fn threshold(corpus_size: u64) -> f64;
pub fn yield_at(corpus_size: u64, effort: f64) -> f64;
\end{lstlisting}

\subsection{\texttt{sebd.rs}}

\begin{lstlisting}[language=rust]
pub struct SEBDResult {
    pub best_costs: Vec<f64>,
    pub meeting_points: Vec<Option<PartitionState>>,
    pub transition_states: Vec<Vec<SEntropy>>,
    pub matched_compounds: Vec<Vec<String>>,
}

pub fn sebd_search(
    precursor: &PartitionState,
    fragments: &[PartitionState],
    trie: &TernaryTrie,
    max_shell_drop: u32,
    trie_depth: usize,
) -> SEBDResult;

pub fn extract_virtual_substates(
    result: &SEBDResult,
) -> Vec<SEntropy>;

pub fn reachability_cone(
    precursor: &PartitionState,
    depth: u32,
) -> Vec<PartitionState>;

pub fn cone_membership(
    cone: &[PartitionState],
    fragments: &[PartitionState],
) -> Vec<bool>;
\end{lstlisting}

% =====================================================================
\part{Validation Architecture}
% =====================================================================

\section{Test Categories}
\label{sec:test-categories}

Every experiment produces three categories of output, each with
distinct validation requirements.

\begin{definition}[Consistency check]
A test that verifies the implementation computes what the formula says.
Passes by construction if the implementation is correct; informative
only when it fails. Examples: Kirchhoff residual $= 0$, relabelling
invariance, step-to-root $=$ address length.
\end{definition}

\begin{definition}[Positive control]
A test on a case where the answer is known independently. Must be
capable of failing on well-formed input. Example: coupon-collector
retrieval is analytically concave, so the diminishing-returns detector
must report concavity there.
\end{definition}

\begin{definition}[Negative control]
A test asserting that something does \emph{not} happen on well-formed
input. Example: the allocator does not refuse above threshold; spreading
does not lose above threshold; the concave branch recovers water-filling
exactly.
\end{definition}

\section{Per-Experiment Validation Design}
\label{sec:per-exp-validation}

\begin{table}[htbp]
\centering
\small
\caption{Validation design per experiment. Each experiment carries
consistency checks, at least one positive control, and at least one
negative control.}
\label{tab:validation-design}
\begin{tabular}{@{}lp{3cm}p{3cm}p{3cm}@{}}
\toprule
\textbf{Experiment} & \textbf{Consistency} & \textbf{Positive control} &
\textbf{Negative control} \\
\midrule
1.\ CASMI &
  InChIKey format; mass accuracy within tolerance &
  Known compounds from NIST must rank 1 &
  Decoy compounds must not rank 1 \\[4pt]
2.\ Trans.\ states &
  Virtual substates average to physical point &
  McLafferty TS geometry matches DFT &
  Random off-shell points do not correlate with DFT \\[4pt]
3.\ Lagrangian &
  $\Sv \in [0,1]^3$ for all compounds &
  Same compound $\rightarrow$ same coordinates on same platform &
  Different compounds $\rightarrow$ different coordinates \\[4pt]
4.\ Glycopeptide &
  Precursor mass conserved in cone &
  Peptide backbone fragments in cone &
  Noise peaks not in cone \\[4pt]
5.\ Scaling &
  Trie depth matches specification &
  Query time flat across DB sizes &
  Empty lookups at off-shell addresses \\[4pt]
6.\ Federated &
  Allocator yield $\leq$ brute force &
  Refusal fires below threshold &
  Refusal does not fire above threshold \\[4pt]
7.\ DSL &
  Programs compile without error &
  Output matches published results &
  Misspelled field names produce warnings \\
\bottomrule
\end{tabular}
\end{table}

\section{Reporting Requirements}
\label{sec:reporting}

Each experiment report must contain, at minimum:

\begin{enumerate}[nosep,label=(\roman*)]
\item The exact version of the Shapeshifter program used, as a
  \shk{.ss} file committed to the repository.
\item The exact version of the Rust core library, as a Git commit hash.
\item The dataset provenance, including download URLs and file checksums.
\item All numerical results as JSON files produced by the validation
  pipeline, not manually transcribed.
\item The comparison method's results, with citation to the source.
\item A table distinguishing consistency checks (pass by construction)
  from controls (capable of failing).
\item Any measurement error or implementation flaw discovered during
  validation, described with the same prominence as positive results
  (following the practice of the attention allocation paper,
  Section~7.3).
\end{enumerate}

% =====================================================================
\part{Implementation Roadmap}
% =====================================================================

\section{Phase 1: Core Infrastructure}
\label{sec:phase1}

\begin{deliverable}[Rust core library: types and ternary addressing]
\label{del:core}
Implement \shk{constants.rs}, \shk{sentropy.rs}, \shk{partition.rs},
\shk{ternary.rs}. These four modules provide the type-level invariants
shared by all subsequent modules. Estimated: 255 lines.
\end{deliverable}

\begin{deliverable}[Trie and lazy dictionary]
\label{del:trie}
Implement \shk{trie.rs} with insert, exact lookup, prefix match, and
empty-lookup detection. Estimated: 100 lines. Required for
Experiments~1 and~5.
\end{deliverable}

\begin{deliverable}[Shapeshifter grammar extensions]
\label{del:grammar}
Add \shk{source}, \shk{network}, and \shk{precursor} block types to
the parser. Add the \shk{foreach} modifier. Update the compile stage
to validate the new blocks. Prove that the four invariant theorems
hold under the extensions.
\end{deliverable}

\section{Phase 2: Experiment-Specific Modules}
\label{sec:phase2}

\begin{deliverable}[SEBD-MS implementation]
\label{del:sebd}
Implement \shk{sebd.rs} and \shk{graph.rs}. The SEBD-MS algorithm,
virtual substate extraction, and forward reachability cone. Required
for Experiments~1, 2, and~4. Estimated: 200 lines.
\end{deliverable}

\begin{deliverable}[Circuit solver and CKG]
\label{del:circuit}
Implement \shk{circuit.rs} and \shk{propagation.rs}. Required for the
CKG-constrained variant of Experiment~1 and the triple-convergence test
of Experiment~3. Estimated: 160 lines.
\end{deliverable}

\begin{deliverable}[Attention allocator]
\label{del:allocator}
Implement \shk{allocator.rs} with gated allocation, refusal, and
brute-force comparison. Required for Experiment~6 and the federated
variant of Experiment~1. Estimated: 80 lines.
\end{deliverable}

\section{Phase 3: Experiments and Publication}
\label{sec:phase3}

\begin{deliverable}[Experiment 2 execution]
\label{del:exp2}
Run the transition state experiment. This is first because it has the
highest novelty-to-effort ratio and requires no external data
collection.
\end{deliverable}

\begin{deliverable}[Experiment 1 execution]
\label{del:exp1}
Run CASMI 2022. This requires Deliverables~\ref{del:trie},
\ref{del:sebd}, and optionally~\ref{del:circuit} and~\ref{del:allocator}.
Produce a spreadsheet in the CASMI submission format (InChIKey or
SMILES per compound) and score against published solutions.
\end{deliverable}

\begin{deliverable}[Results manuscript]
\label{del:manuscript}
Write a single results paper covering Experiments~1--6, structured as:
introduction (the gap), methods (the framework, briefly, with citations
to the theory papers), results (per-experiment, with comparison tables),
and discussion (what works, what fails, and what each failure localises).
Target venue: Analytical Chemistry or the Journal of Chemical Information
and Modeling.
\end{deliverable}

% =====================================================================
\section*{Summary of All Specifications}
% =====================================================================

\begin{table}[htbp]
\centering
\small
\caption{Complete specification summary.}
\label{tab:summary}
\begin{tabular}{@{}lr@{}}
\toprule
\textbf{Component} & \textbf{Count} \\
\midrule
Experiments defined & 7 \\
Success criteria stated & 7 \\
Comparison methods named & 6 (Exp.~2 has none: novel) \\
DSL block types added & 3 (\shk{source}, \shk{network}, \shk{precursor}) \\
DSL structural extensions & 1 (\shk{foreach} modifier) \\
Failure categories & 3 (warning, error, refusal) \\
Standard library namespaces & 12 (7 existing + 5 new) \\
Standard library operations & 61 (13 existing + 48 new) \\
Rust modules & 14 \\
Estimated Rust lines & 895 \\
Invariant theorems preserved & 4 of 4 \\
Invariant theorems modified & 1 (phase isolation, scoped exception) \\
\bottomrule
\end{tabular}
\end{table}

\bibliographystyle{plainnat}
\bibliography{references}

\end{document}